\section{Locally noetherian formal schemes}
\source{fga-explained:page-0191}\source{fga-explained:page-0192}\source{fga-explained:page-0193}\source{fga-explained:page-0194}\source{fga-explained:page-0195}\source{fga-explained:page-0196}\source{fga-explained:page-0197}\source{fga-explained:page-0198}\source{fga-explained:page-0199}\source{fga-explained:page-0200}\source{fga-explained:page-0201}\source{fga-explained:page-0202}\source{fga-explained:page-0203}\source{fga-explained:page-0204}\source{fga-explained:page-0205}\source{fga-explained:page-0206}\source{fga-explained:page-0207}\source{fga-explained:page-0208}\source{fga-explained:page-0209}\source{fga-explained:page-0210}\source{fga-explained:page-0211}\source{fga-explained:page-0212}\source{fga-explained:page-0213}\source{fga-explained:page-0214}\source{fga-explained:page-0215}\source{fga-explained:page-0216}\source{fga-explained:page-0217}\source{fga-explained:page-0218}\source{fga-explained:page-0219}\source{fga-explained:page-0220}\source{fga-explained:page-0221}\source{fga-explained:page-0222}\source{fga-explained:page-0223}\source{fga-explained:page-0224}\source{fga-explained:page-0225}\source{fga-explained:page-0226}\source{fga-explained:page-0227}\source{fga-explained:page-0228}\source{fga-explained:page-0229}\source{fga-explained:page-0230}\source{fga-explained:page-0231}\source{fga-explained:page-0232}\source{fga-explained:page-0233}\source{fga-explained:page-0234}\source{fga-explained:page-0235}\source{fga-explained:page-0236}\source{fga-explained:page-0237}\source{fga-explained:page-0238}\source{fga-explained:page-0239}\source{fga-explained:page-0240}\source{fga-explained:page-0241}\source{fga-explained:page-0242}\source{fga-explained:page-0243}\source{fga-explained:page-0244}\source{fga-explained:page-0245}\source{fga-explained:page-0246}

\begin{definition}[adic noetherian ring]\label{fga-def-8-1-1}\source{fga-explained:page-0191}\uses{}
An adic noetherian ring is a noetherian ring $A$ complete and separated for the
$I$-adic topology of an ideal of definition $I$; thus
$A\simeq\varprojlim_n A/I^{n+1}$.  A locally noetherian formal scheme is locally
$\operatorname{Spf}(A)$ for such a ring.
\end{definition}\begin{definition}[formal completion]\label{fga-def-8-1-2}\source{fga-explained:page-0192}\uses{fga-def-8-1-1}
For a noetherian scheme $X$ and closed subscheme $Y$ defined by $I$, its formal
completion is $\widehat X=\varprojlim_n(X,\OO_X/I^{n+1})$ with thickenings
$X_n=(X,\OO_X/I^{n+1})$.
\end{definition}

\section{The comparison theorem}
\begin{theorem}[formal functions and comparison]\label{fga-thm-8-2-2}\dcref{8.2.2}\source{fga-explained:page-0198}\uses{fga-def-8-1-2,fga-thm-5-6}
Let $f:X\to Y$ be a finite type morphism of noetherian schemes and $Y'
\hookrightarrow Y$ a closed subscheme.  For a coherent $\mathcal F$ on $X$, the
$(I)$-adic completion of $R^if_*\mathcal F$ identifies with
$\varprojlim_nR^if_{n*}(\mathcal F/I^{n+1}\mathcal F)$ when $f$ is proper.
\end{theorem}
\begin{proof}
Resolve $\mathcal F$ by finite sums of relatively negative twists and use
Serre vanishing to compute each cohomology group by a finite complex.  Completion
commutes with kernels and cokernels in the noetherian category, and passage to
the inverse limit is exact by Mittag--Leffler, giving the claimed isomorphism.
\end{proof}
\begin{corollary}\label{fga-cor-8-2-4}\dcref{8.2.4}\source{fga-explained:page-0200}\uses{fga-thm-8-2-2}
For a proper morphism over a complete local noetherian base, coherent sheaf
cohomology is the inverse limit of the cohomologies on the infinitesimal
thickenings.
\end{corollary}
\begin{proposition}\label{fga-prop-8-2-5-3}\dcref{8.2.5.3}\source{fga-explained:page-0201}\uses{fga-def-8-1-1}
An inverse system of coherent sheaves with surjective transition maps on a
noetherian formal scheme has coherent inverse-limit module on affine opens.
\end{proposition}
\begin{theorem}\label{fga-thm-8-2-6-1}\dcref{8.2.6.1}\source{fga-explained:page-0202}\uses{fga-thm-8-2-2,fga-prop-8-2-5-3}
If $f:X\to Y$ is proper with $Y$ noetherian, completion induces an equivalence
between coherent sheaves on $X$ supported formally along a closed subscheme and
coherent sheaves on the formal completion $\widehat X$.
\end{theorem}
\begin{corollary}\label{fga-cor-8-2-9}\dcref{8.2.9}\source{fga-explained:page-0204}\uses{fga-thm-8-2-6-1}
The natural map from algebraic to formal global sections is dense and induces
the comparison isomorphism on completed modules.
\end{corollary}
\begin{corollary}\label{fga-cor-8-2-10}\dcref{8.2.10}\source{fga-explained:page-0204}\uses{fga-cor-8-2-9}
The theorem on formal functions identifies completed direct images with formal
direct images and is compatible with arbitrary proper base change.
\end{corollary}
\begin{theorem}[Zariski connectedness]\label{fga-thm-8-2-12}\dcref{8.2.12}\source{fga-explained:page-0205}\uses{fga-thm-8-2-2}
If $f:X\to Y$ is proper with geometrically connected fibers and $Y$ is normal,
then the formal and algebraic connectedness of the fibers agree; in particular
the Stein factor has geometrically connected fibers.
\end{theorem}
\begin{corollary}\label{fga-cor-8-2-17}\dcref{8.2.17}\source{fga-explained:page-0207}\uses{fga-thm-8-2-12}
Zariski's main theorem holds in the form that a proper quasi-finite morphism is
finite after replacing the target by the Stein factor.
\end{corollary}
\begin{corollary}\label{fga-cor-8-2-18}\dcref{8.2.18}\source{fga-explained:page-0208}\uses{fga-thm-8-2-12}
Over a henselian noetherian local base, finite etale covers of the closed fiber
extend uniquely to finite etale covers of the scheme.
\end{corollary}

\section{Cohomological flatness}
\begin{theorem}[cohomological flatness]\label{fga-thm-8-3-2}\dcref{8.3.2}\source{fga-explained:page-0209}\uses{fga-thm-5-7}
For a proper morphism $f:X\to S$ and coherent $\mathcal F$ flat over $S$, the
following are equivalent near $s\in S$: $f_*\mathcal F$ is locally free and
commutes with base change; the fiberwise $H^0$ dimension is locally constant;
and $\mathcal F$ is cohomologically flat in dimension zero.
\end{theorem}
\begin{corollary}\label{fga-cor-8-3-3}\dcref{8.3.3}\source{fga-explained:page-0209}\uses{fga-thm-8-3-2}
Cohomological flatness in dimension zero is stable under arbitrary base change.
\end{corollary}
\begin{corollary}\label{fga-cor-8-3-4}\dcref{8.3.4}\source{fga-explained:page-0210}\uses{fga-cor-8-3-3}
If $f$ is proper and flat and $\mathcal F$ is invertible, cohomological
flatness in dimension zero is equivalent to the fibers having constant number
of global sections.
\end{corollary}
\begin{proposition}\label{fga-prop-8-3-6-4}\dcref{8.3.6.4}\source{fga-explained:page-0211}\uses{fga-def-8-1-1}
For a pseudo-coherent complex $E$, perfectness at a point is equivalent to
finite Tor-amplitude and finite-dimensional derived fibers there.
\end{proposition}
\begin{corollary}\label{fga-cor-8-3-6-5}\dcref{8.3.6.5}\source{fga-explained:page-0211}\uses{fga-prop-8-3-6-4}
Perfect complexes are stable under tensor product, duality, and proper derived
pushforward in the bounded coherent setting.
\end{corollary}
\begin{theorem}\label{fga-thm-8-3-8}\dcref{8.3.8}\source{fga-explained:page-0213}\uses{fga-cor-8-3-6-5,fga-thm-8-2-2}
For a proper morphism, a perfect complex has a bounded perfect derived direct
image locally on the base, and its formation commutes with derived base change.
\end{theorem}
\begin{corollary}\label{fga-cor-8-3-11}\dcref{8.3.11}\source{fga-explained:page-0215}\uses{fga-thm-8-3-8}
Under cohomological flatness, the derived fibers of a proper perfect complex
have constant Euler characteristic and the base-change maps are isomorphisms.
\end{corollary}

\section{The existence theorem}
\begin{theorem}[algebraization of formal coherent sheaves]\label{fga-thm-8-4-2}\dcref{8.4.2}\source{fga-explained:page-0218}\uses{fga-thm-8-2-6-1,fga-cor-8-3-11}
Let $X$ be a proper scheme over a complete noetherian local ring and
$\widehat X$ its formal completion.  Completion induces an equivalence between
coherent sheaves on $X$ and coherent sheaves on $\widehat X$.
\end{theorem}
\begin{proof}
For a formal coherent sheaf $\mathcal F=(\mathcal F_n)$, choose a sufficiently
positive relatively ample twist so that all $\mathcal F_n(m)$ are generated
and have vanishing higher cohomology.  The comparison theorem makes the spaces
of sections a compatible inverse system; a finite presentation descends from a
single level and the relations descend by completeness.  The resulting coherent
algebraic sheaf completes to $\mathcal F$, and full faithfulness follows from
formal functions applied to internal Hom.
\end{proof}
\begin{lemma}\label{fga-lem-8-4-3-1}\dcref{8.4.3.1}\source{fga-explained:page-0217}\uses{fga-thm-8-2-2}
For a coherent formal sheaf, there is a uniform twist making all levels globally
generated with vanishing higher cohomology; the multiplication maps are
surjective compatibly with the inverse system.
\end{lemma}
\begin{corollary}\label{fga-cor-8-4-5}\dcref{8.4.5}\source{fga-explained:page-0221}\uses{fga-thm-8-4-2}
Algebraization is fully faithful on coherent sheaves and on their finite
presentations.
\end{corollary}
\begin{corollary}\label{fga-cor-8-4-6}\dcref{8.4.6}\source{fga-explained:page-0221}\uses{fga-cor-8-4-5}
Invertible formal sheaves and their sections algebraize uniquely.
\end{corollary}
\begin{corollary}\label{fga-cor-8-4-7}\dcref{8.4.7}\source{fga-explained:page-0222}\uses{fga-cor-8-4-6}
Proper formal schemes over a complete noetherian base with an ample invertible
sheaf are algebraizable as proper schemes.
\end{corollary}
\begin{theorem}\label{fga-thm-8-4-10}\dcref{8.4.10}\source{fga-explained:page-0224}\uses{fga-thm-8-4-2}
The algebraization equivalence is compatible with morphisms, closed immersions,
and completion of proper schemes over the base.
\end{theorem}

\section{Applications to lifting problems}
\begin{theorem}\label{fga-thm-8-5-3}\dcref{8.5.3}\source{fga-explained:page-0226}\uses{fga-thm-8-4-2}
For a proper scheme over a complete local base, a compatible formal system of
deformations together with a compatible ample invertible sheaf algebraizes to a
unique algebraic deformation.
\end{theorem}
\begin{corollary}\label{fga-cor-8-5-5}\dcref{8.5.5}\source{fga-explained:page-0227}\uses{fga-thm-8-5-3}
Formal liftings of proper schemes over complete local rings are algebraizable
when the special fiber has a formal ample line bundle.
\end{corollary}
\begin{theorem}\label{fga-thm-8-5-12}\dcref{8.5.12}\source{fga-explained:page-0230}\uses{fga-def-6-4-2,fga-thm-8-4-2}
For a smooth proper scheme, the obstruction to lifting across a square-zero
extension lies in $\operatorname{Ext}^2(L_{X/S},J)$; if it vanishes, liftings
form a torsor under $\operatorname{Ext}^1(L_{X/S},J)$.
\end{theorem}
\begin{theorem}\label{fga-thm-8-5-15}\dcref{8.5.15}\source{fga-explained:page-0231}\uses{fga-thm-8-2-12,fga-thm-8-4-10}
Grothendieck's specialization theorem identifies formal and algebraic
specializations of proper smooth families and preserves connected components.
\end{theorem}
\begin{proposition}\label{fga-prop-8-5-16}\dcref{8.5.16}\source{fga-explained:page-0232}\uses{fga-thm-8-5-12}
For proper flat $X\to Y$, a line bundle with fiberwise constant Hilbert
polynomial extends uniquely over infinitesimal thickenings when its obstruction
class vanishes.
\end{proposition}
\begin{theorem}\label{fga-thm-8-5-19}\dcref{8.5.19}\source{fga-explained:page-0234}\uses{fga-thm-8-5-3,fga-thm-8-5-12}
A proper smooth family over a complete local ring admits algebraic liftings of
its formal family; for curves the obstruction is controlled by the Jacobian and
vanishes after the stated cohomological hypotheses.
\end{theorem}
\begin{theorem}\label{fga-thm-8-5-20}\dcref{8.5.20}\source{fga-explained:page-0235}\uses{fga-thm-8-5-19}
The formal lifting criterion applies to smooth projective varieties with an
ample line bundle and gives a projective algebraic lifting over the complete
base whenever the obstruction groups vanish.
\end{theorem}
\begin{theorem}\label{fga-thm-8-5-27}\dcref{8.5.27}\source{fga-explained:page-0237}\uses{fga-thm-8-5-12}
For the cotangent complex, the obstruction to lifting a morphism is the class in
$\operatorname{Ext}^2(L_{X/S},J)$ and the set of liftings is a torsor under
$\operatorname{Ext}^1$; this extends the smooth case.
\end{theorem}
\begin{corollary}\label{fga-cor-8-5-32}\dcref{8.5.32}\source{fga-explained:page-0238}\uses{fga-thm-8-5-27,fga-thm-8-4-2}
For a proper smooth curve, formal deformations with fixed special fiber and
polarization algebraize; the completed cotangent complex agrees with the
algebraic one.
\end{corollary}

\section{Serre's examples and letter}
\begin{proposition}\label{fga-prop-8-6-2}\dcref{8.6.2}\source{fga-explained:page-0239}\uses{fga-thm-8-5-12}
The Raynaud--Serre construction gives smooth projective varieties in
characteristic $p>0$ with nonzero obstruction to lifting to characteristic zero.
\end{proposition}
\begin{proposition}\label{fga-prop-8-6-4}\dcref{8.6.4}\source{fga-explained:page-0239}\uses{fga-cor-8-4-7}
Under the numerical hypotheses on $r,p,n$ in the source, a formal lifting over a
complete local noetherian ring is algebraizable and projective.
\end{proposition}
\begin{proposition}\label{fga-prop-8-6-6}\dcref{8.6.6}\source{fga-explained:page-0240}\uses{fga-prop-8-6-2}
For $p>n+1$, the specified cyclic covers and their canonical bundles yield the
non-liftable examples; the obstruction is detected by the Hodge filtration.
\end{proposition}
\begin{corollary}\label{fga-cor-8-6-7}\dcref{8.6.7}\source{fga-explained:page-0241}\uses{fga-prop-8-6-6}
The resulting smooth projective varieties are not liftable to characteristic
zero despite their formal local liftings.
\end{corollary}
